
\begin{obereflection}
\textbf{Latihan 11.1}
1. Mengapa pemilihan akar primitif $g$ sangat penting dalam protokol Diffie-Hellman?
2. Jelaskan perbedaan mendasar antara peran kunci publik dan kunci privat dalam RSA dibandingkan dengan dalam protokol Diffie-Hellman.
3. Bagaimana cara mencegah serangan \textit{Man-in-the-Middle} pada pertukaran kunci Diffie-Hellman? (Petunjuk: Pikirkan tentang tanda tangan digital atau sertifikat).
4. Sebuah organisasi memiliki 500 pegawai yang perlu berkomunikasi rahasia secara
   berpasangan memakai AES. Hitunglah berapa kunci rahasia yang harus dikelola bila
   setiap pasangan memakai kunci tersendiri. Bandingkan dengan jumlah kunci bila
   organisasi itu memakai pasangan kunci Diffie--Hellman sesaat untuk setiap sesi.
5. Jelaskan mengapa memperbesar modulus $p$ tidak menghalangi serangan
   \textit{Man-in-the-Middle}, padahal memperbesar $p$ menghalangi penyadapan pasif.
   Apa yang sesungguhnya menjadi celah keamanan pada serangan tersebut?
6. Pada protokol Diffie--Hellman sesaat (DHE), kunci privat $a$ dan $b$ dibuang
   setelah sesi berakhir. Jelaskan bagaimana kebiasaan itu melindungi rekaman
   percakapan yang sudah lama tersimpan pada saat kunci jangka panjang peladen
   akhirnya bocor.
\end{obereflection}

\begin{obereflection}
\textbf{Latihan 11.2}
1. Perhatikan $p = 23$, yang merupakan \textit{safe prime} karena $23 = 2 \times 11 + 1$
   dengan $q = 11$ juga prima. Tunjukkan bahwa $g = 2$ bukan akar primitif dari $23$,
   lalu hitunglah ordenya. Jelaskan mengapa $g$ semacam itu justru dianjurkan pada
   protokol yang memakai subgrup berorde prima.
2. Diketahui $p = 23$ dan $g = 5$. Tentukan seluruh akar primitif dari $23$ dan
   periksa bahwa jumlahnya sama dengan $\varphi(22)$. Tunjukkan bahwa $5$ termasuk
   di dalamnya.
3. Diketahui $p = 23$, $g = 5$, $a = 7$, dan $b = 15$. Hitunglah $A$, $B$, dan kunci
   bersama $K$ dari kedua sisi. Periksa hasil Anda dengan
   Contoh~\ref{ex:dh-kecil} dan dengan program pada Listing~\ref{lst:diffie-hellman}.
4. Modulus $p = 353$ memiliki $160$ akar primitif, yaitu sebanyak $\varphi(352)$.
   Jelaskan mengapa jumlah akar primitif modulo $p$ selalu sama dengan
   $\varphi(p-1)$, dan apa akibat praktisnya bagi pihak yang memilih $g$ secara acak.
5. Dengan algoritma langkah raksasa-bayi, penyerang memerlukan sekitar
   $2\sqrt{p}$ langkah. Untuk modulus 2048 bit, tunjukkan bahwa $\sqrt{2^{2048}}$
   berukuran sekitar $1{,}8 \times 10^{308}$. Bila satu komputer dapat melakukan
   $10^{9}$ perkalian per detik, perkirakan berapa tahun waktu yang diperlukan.
   Bandingkan dengan umur alam semesta yang sekitar $1{,}4 \times 10^{10}$ tahun.
6. Pada contoh $p = 353$, penyerang menemukan $a = 97$ hanya dengan $24$
   perkalian langkah raksasa-bayi, jauh lebih sedikit daripada $97$ langkah yang
   diperlukan pencarian menyeluruh. Bila Alice memilih kunci privat yang lebih
   besar, misalnya $a = 300$, jumlah perkaliannya memang bertambah, tetapi tetap
   berorde $\sqrt{p}$ dan bukan berorde $a$. Jelaskan mengapa demikian, dan apa
   yang sebenarnya diukur oleh perkiraan $2\sqrt{p}$ itu.
7. Rancanglah satu skema sederhana yang memungkinkan Alice dan Bob memastikan
   keaslian nilai pertukaran mereka dengan memakai tanda tangan digital. Jelaskan
   bagian mana dari skema Anda yang membuat Mallory tidak dapat menggantikan
   $g^{a}$ dengan $g^{m}$.
\end{obereflection}
